\section{Discussion, limitations and conclusion}
\label{sec:discussion}

\subsection{What generalises}

Three lessons seem portable beyond this application.

\paragraph{An incremental format turns every writer into a stateful operation.} The bug in
\S\ref{sec:incremental} is not a coding error in the ordinary sense; it is a failure to model
the writer as having a \emph{past}. A function that appends is correct once, so any provider of
such a function must answer the question ``what did you write last time?'' The answer in this
system is a per-page set of object numbers held on the document session, and the general
statement is that the state belongs next to the bytes, not inside the writer.

\paragraph{Flattening a tree materialises inheritance.} This is true of any format with
inheritable attributes --- CSS, DOM, scene graphs, PDF page trees. The failure is silent
because the flattened document is still well-formed; the flaw is that it means something else.

\paragraph{A test can be wrong in a way that looks like a finding.} The projection trap in
\S\ref{sec:verification} cost the project more time than any product defect, twice, on the same
feature pattern. The rule we now apply --- assert against the model, not against a debug
surface that describes the model --- is cheap and should have been adopted on the first
occurrence.

\subsection{Limitations}
\label{sec:limits}

This list is the boundary of what the paper's claims cover.

\begin{itemize}
  \item \textbf{No installer and no code signing.} The application ships as a portable
        executable; Windows SmartScreen warns on first run. Fixing this needs a code-signing
        certificate, not code.
  \item \textbf{Windows only.} The packaged target is Windows x64. Nothing in the renderer is
        platform-specific, but the shell is not exercised elsewhere.
  \item \textbf{Encrypted documents are refused for writing.} They open for reading through the
        renderer; an annotation save against an encrypted file is rejected rather than
        attempted.
  \item \textbf{Files grow with every save.} The incremental writer never rewrites the original
        bytes, so superseded objects remain as orphans. A compaction pass would need a
        whole-document rewrite and is not implemented.
  \item \textbf{Flattening cannot encode text outside Latin-1.} The embedded appearance font is
        a standard Helvetica; a flatten of Chinese or Greek annotation text renders as
        placeholders. The summary feature avoids this by rendering pages instead of writing
        text (\S\ref{sec:pdf-layer}).
  \item \textbf{No form filling, digital signatures or OCR.} Form fields are highlighted, and
        \S\ref{sec:fields} is explicit about why that is not the same as supporting forms.
  \item \textbf{Replace text is written as \texttt{/StrikeOut}}. The format has no replace-text
        annotation type; another reader shows the strike and not the substitution, which is
        visible in the PDF and correct, and incomplete as a representation of the intent.
  \item \textbf{One machine, one engine.} The timings come from a single Windows host and a
        single pdf.js version.
  \item \textbf{Rendering quality is not evaluated.} The application delegates rasterisation to
        pdf.js; any rendering difference the reader observes is the engine's, not this
        project's, and the paper makes no claim either way.
\end{itemize}

\subsection{Conclusion}

A PDF reader that writes is a different system from a PDF reader that renders, and the
difference is concentrated in a small number of invariants that are easy to state and easy to
break: rebuild a page's annotation array rather than appending to it; materialise what a page
inherits before flattening its tree; renumber references when copying between documents and
leave stream bytes alone. Each invariant, violated, produces a file that opens and is wrong ---
which is why the verification described here asks questions of pixels and bytes rather than of
return values, and why the paper reports the defects that reached a user as well as the ones
the suites caught.

The application is 17{,}323 lines of source over 47 files with one runtime dependency, MIT
licensed, and available with its tests and the scripts that rebuilt every figure here at \repo.

\section*{Acknowledgements}

The software described here was developed with the assistance of \textbf{DeepSeek AI} and of
\textbf{DSH} (DeepSeek Harness), the agent harness built on it, and this preprint was
researched, drafted and revised with the same assistance. The debt is not incidental to the
result. The invariant that makes annotation deletion work (\S\ref{sec:incremental}), the
decision to materialise inheritance before flattening a page tree (\S\ref{sec:rewrite}), and
the harness defects recorded in \S\ref{sec:verification} were all arrived at in the course of
that collaboration, and the manuscript was argued out the same way.

It is in that spirit --- and because a program that expects its behaviour to be checked should
let its readers check it --- that all of this software is released as open source, with nothing
held back. The disclosure is deliberate: a reader is entitled to know how a document was
produced, as a reader of the source is entitled to see every line of it.

The dual representation of \S\ref{sec:annotations}, the coordinate convention and the decision
to write into the document rather than beside it came from using the software daily while
building it. The pdf.js project~\cite{pdfjs} does the rendering; the writer's invariants were
learned by breaking them.

I thank Liang Shengren, and this remarkable era, in which such an undertaking is within
reach of one person with a laptop and a good interlocutor.

